\begin{frame}{1969, \textit{Perceptrons}: XOR의 벽과 AI의 겨울}
  \begin{itemize}
    \item \kb{민스키·페퍼트(MIT, 1969)}: 단층 퍼셉트론은 XOR을
      \textbf{아무리 학습시켜도 절대 못 맞춘다} — 수학적으로 증명
    \item ``학습률 부족'', ``데이터 부족''이 아니라 \textbf{모델 구조 자체}가 불가능
  \end{itemize}
  \vskip0.5em
  \kq{영향: 투자·논문 급감 $\to$ 1차 ``AI 겨울''(1970s 후반$\sim$1980s 중반)의 직접 원인.
  대중적으로는 ``신경망은 쓸모없다''로 과도하게 일반화됐다.}
  \vskip0.5em
  \begin{itemize}
    \item 그 긴 동면기를 깨운 것: 1986년 \kb{럼멜하트-힌튼-윌리엄스}의
      \kb{역전파} 대중화 논문
    \item 동면기를 깨운 것이 바로 ``XOR의 벽을 넘는 단 하나의 구조'' —
      다음 프레임의 \kb{은닉층}
  \end{itemize}
\end{frame}
